\documentclass[a4paper,10pt,brazil]{article}
\usepackage{provastyle}
\begin{document}
\makeexamtitlepage{UFOP}{BCC222: Programação Funcional}{José Romildo}{2025/2}{Exame Especial Parcial 1}{03}
\begin{multicols}{2}
\questionTitle{1}{3: Expressões e Definições}
Haskell possui regras sintáticas rígidas para diferenciar nomes (identificadores) dados a variáveis, funções, tipos e construtores. De acordo com as regras para identificadores \textbf{alfanuméricos}, qual das alternativas apresenta um nome \textbf{válido e correto} para definir uma nova \textbf{função} local ou de nível superior?

\begin{enumerate}
  \item \haskellinline{1oCalculo}

  \item \haskellinline{_calculoFinal'}

  \item \haskellinline{calculo-final}

  \item \haskellinline{:calculoFinal}
  \item \haskellinline{CalculoFinal}


\end{enumerate}

\questionTitle{2}{9: Recursividade}
Uma das propriedades mais poderosas do Haskell é sua interação entre chamadas recursivas e a avaliação preguiçosa (\emph{lazy evaluation}). Considere a função a seguir, que \textbf{não possui um caso base}:

\begin{minted}{haskell}
geraSequencia :: Int -> [Int]
geraSequencia n = n : geraSequencia (n + 2)
\end{minted}

O que acontecerá se você solicitar ao compilador para avaliar a expressão \haskellinline{take 3 (geraSequencia 10)}?

\begin{enumerate}
  \item Ocorrerá um erro de execução na terceira chamada, pois o Haskell impõe um limite de segurança de no máximo duas chamadas recursivas para funções que não declaram caso base explícito.
  \item O programa travará em um laço infinito imediato (\emph{Stack overflow}) e consumirá toda a memória antes de repassar o resultado para a função \haskellinline{take}.

  \item O resultado será \haskellinline{[10, 12, 14]}. A ausência do caso base não causa um laço infinito, pois a avaliação preguiçosa expandirá a recursão apenas o número de vezes estritamente necessário para que a função \haskellinline{take 3} seja satisfeita.

  \item O compilador rejeitará a função \haskellinline{geraSequencia} em tempo de compilação emitindo um aviso de \emph{Missing base case for recursive function}.

  \item O resultado será \haskellinline{[10, 10, 10]}, porque a função atualiza o valor de \haskellinline{n} no escopo local, mas não no escopo da chamada geradora original.


\end{enumerate}

\questionTitle{3}{10: Ações de E/S Recursivas}
O texto base enfatiza a importância de separar o código de E/S do código de processamento puro. Suponha que você precise ler várias linhas numéricas digitadas pelo usuário, ordená-las com o algoritmo Quicksort e imprimir o resultado. 

Por que a melhor prática em Haskell dita que a ordenação (Quicksort) \textbf{não} deve ser misturada dentro das equações da ação recursiva de leitura?

\begin{enumerate}
  \item Não há benefício técnico prático além de formatação visual; em tempo de compilação, o Haskell funde as funções puras e de E/S no mesmo laço imperativo de máquina.

  \item Porque o compilador GHC proíbe chamadas de funções matemáticas dentro de blocos \haskellinline{do}; apenas funções que retornam o tipo \haskellinline{IO} são permitidas sintaticamente.

  \item Para isolar os efeitos colaterais. Mantendo o Quicksort como uma função puramente funcional (ex: \haskellinline{[Int] -> [Int]}), ela permanece matematicamente previsível, fácil de testar independentemente e imune a falhas de entrada do teclado.

  \item Porque o Quicksort consome muita memória na pilha e, se for chamado recursivamente junto com operações de E/S, causará invariavelmente um \emph{overflow} no buffer de leitura do sistema operacional.

  \item Para evitar a conversão excessiva de tipos. Se a ordenação ficasse na função de leitura, ela seria obrigada a operar sobre listas do tipo \haskellinline{[IO Int]} em vez de \haskellinline{[Int]}.


\end{enumerate}

\questionTitle{4}{1: Paradigmas}
Dada a definição da função \haskellinline{soma} abaixo, qual é a sequência de redução passo a passo logicamente correta para provar que a expressão \haskellinline{soma [5]} resulta em \haskellinline{5}?

\begin{minted}{haskell}
soma []     = 0
soma (x:xs) = x + soma xs
\end{minted}

\begin{enumerate}
  \item \haskellinline{soma (5:0)} $\leadsto$ \haskellinline{5 + soma 0} $\leadsto$ \haskellinline{5 + []} $\leadsto$ \haskellinline{5}

  \item \haskellinline{soma [5, []]} $\leadsto$ \haskellinline{5 + []} $\leadsto$ \haskellinline{5 + 0} $\leadsto$ \haskellinline{5}

  \item \haskellinline{soma (5:[])} $\leadsto$ \haskellinline{5 + soma []} $\leadsto$ \haskellinline{5 + 0} $\leadsto$ \haskellinline{5}

  \item \haskellinline{soma [5]} $\leadsto$ \haskellinline{5 + 0} $\leadsto$ \haskellinline{soma []} $\leadsto$ \haskellinline{5}

  \item \haskellinline{soma [5]} $\leadsto$ \haskellinline{5 + soma 5} $\leadsto$ \haskellinline{5 + 0} $\leadsto$ \haskellinline{5}


\end{enumerate}

\questionTitle{5}{7: Polimorfismo}
Um aluno tentou escrever uma função para calcular a média de uma lista de números fracionários:
\begin{minted}{haskell}
media :: Fractional a => [a] -> a
media xs = sum xs / length xs
\end{minted}
O compilador rejeita este código com um erro de tipo. Qual é a causa exata do erro e como resolvê-lo mantendo a precisão fracionária?

\begin{enumerate}
  \item O operador (\haskellinline{/}) não é polimórfico. Ele só funciona exclusivamente para o tipo \haskellinline{Double}. A solução é aplicar a função \haskellinline{realToFrac} na variável \haskellinline{xs} inteira.

  \item A função \haskellinline{length} pode retornar 0 para listas vazias, e o compilador prevê a divisão por zero. A solução é adicionar uma guarda que trate o caso de lista vazia.

  \item A função \haskellinline{length} retorna um tipo estrito \haskellinline{Int}, mas o operador de divisão (\haskellinline{/}) exige que ambos os operandos pertençam à classe \haskellinline{Fractional}. A solução é usar \haskellinline{fromIntegral (length xs)}.

  \item A função \haskellinline{sum} retorna um número inteiro por padrão. O operador (\haskellinline{/}) falha porque tenta dividir um inteiro por outro. A solução é alterar o tipo da função para \haskellinline{media :: [Int] -> Double}.

  \item O erro ocorre porque listas em Haskell não podem ser usadas em operações matemáticas diretas. A solução é usar um laço de repetição \haskellinline{for} para somar os valores individualmente.


\end{enumerate}

\questionTitle{6}{3: Expressões e Definições}
Considere a definição de uma função \haskellinline{f} em Haskell que contém múltiplas definições locais na sua cláusula \haskellinline{where}. 
\begin{minted}{haskell}
f x = a + b
  where
    b = a * 2
    a = x + 10
\end{minted}
Sobre o sistema de resolução de dependências de Haskell para o bloco \haskellinline{where}, é correto afirmar que:

\begin{enumerate}
  \item Ocorrerá um erro de compilação (\emph{Not in scope}), pois a variável \haskellinline{a} é referenciada na definição de \haskellinline{b} antes de ter sido declarada.

  \item A cláusula \haskellinline{where} só permite a definição de uma única variável local. Para múltiplas definições, deve-se usar blocos \haskellinline{let ... in} aninhados.

  \item O código causará um laço infinito em tempo de execução devido à recursividade mútua não resolvida entre as variáveis \haskellinline{a} e \haskellinline{b}.

  \item A ordem das definições locais não importa; Haskell resolve o sistema de equações e dependências automaticamente, permitindo que a expressão seja avaliada sem erros.

  \item O resultado será corrompido, pois \haskellinline{a} receberá um valor nulo por padrão quando avaliado fora de ordem sintática.


\end{enumerate}

\questionTitle{7}{2: Ambiente}
Durante o desenvolvimento de um projeto, o seu editor de código (Visual Studio Code) parou subitamente de exibir marcações de erro de sintaxe em tempo real e perdeu a funcionalidade de autocompletar o código. No entanto, você verifica que o projeto ainda compila normalmente quando você executa \texttt{cabal build} diretamente no terminal. 

Qual ferramenta do ecossistema Haskell, instalada pelo GHCup, provavelmente falhou ou foi desativada no seu editor?

\begin{enumerate}
  \item O compilador GHC, que é o único componente capaz de analisar o código-fonte em tempo de edição para fornecer avisos.

  \item O GHCi, pois é ele que roda em segundo plano na IDE para verificar erros de sintaxe e inferir tipos durante a digitação.

  \item O Cabal, pois é a ferramenta que gerencia as dependências visuais e a interface com as extensões do Visual Studio Code.

  \item O GHCup, que precisa estar em execução contínua como um serviço (daemon) para manter a IDE sincronizada com o compilador.

  \item O HLS (Haskell Language Server), que é o serviço responsável por fornecer inteligência de código e integração dinâmica com a IDE.


\end{enumerate}

\questionTitle{8}{7: Polimorfismo}
Em Haskell, os literais numéricos no código não têm um tipo rígido predefinido, sendo sobrecarregados através de classes de tipo. O literal \haskellinline{10} tem o tipo inferido \haskellinline{Num a => a}, enquanto o literal \haskellinline{10.5} tem o tipo \haskellinline{Fractional a => a}. 
O que acontecerá se você forçar o tipo escrevendo a expressão \haskellinline{10.5 :: Int}?

\begin{enumerate}
  \item A expressão será avaliada como \haskellinline{11}, pois a conversão explícita usando \haskellinline{::} força o arredondamento simétrico do valor fracionário para o inteiro mais próximo.

  \item Ocorrerá um erro de tipo em tempo de compilação, pois o tipo \haskellinline{Int} pertence à classe \haskellinline{Num}, mas não é uma instância da classe \haskellinline{Fractional}, exigida pela presença do ponto decimal.

  \item A expressão será avaliada como \haskellinline{10}, pois o Haskell converterá automaticamente o valor truncando a parte decimal, simulando o comportamento da linguagem C.

  \item Ocorrerá um erro em tempo de execução, já que a linguagem tentará procurar a implementação da função \haskellinline{truncate} na classe \haskellinline{Int} e não a encontrará.

  \item A expressão será aceita sem erros, mas o valor interno será armazenado como um número de ponto flutuante, invalidando a promessa da tipagem forte até que seja exibido na tela.


\end{enumerate}

\questionTitle{9}{6: Estruturas de dados básicas}
Você está modelando um sistema escolar. Definiu o tipo \haskellinline{Aluno} como uma tupla contendo o nome, uma lista com as notas dos testes, e a nota do exame final, que é opcional (o aluno pode não ter precisado fazer):
\begin{minted}{haskell}
type Aluno = (String, [Double], Maybe Double)
\end{minted}

Qual das seguintes expressões cria uma representação válida e tipagem correta de um \haskellinline{Aluno} chamado \haskellinline{"Carlos"}, que tirou 7.0 e 8.5 nos testes, e \textbf{não precisou} fazer o exame final?

\begin{enumerate}
  \item \haskellinline{["Carlos", 7.0, 8.5, Nothing]}

  \item \haskellinline{("Carlos", [7.0, 8.5])}

  \item \haskellinline{("Carlos", [7.0, 8.5], Nothing)}

  \item \haskellinline{("Carlos", [7.0, 8.5], Just 0.0)}

  \item \haskellinline{("Carlos", [7.0, 8.5], Maybe)}


\end{enumerate}

\questionTitle{10}{4: Tipos de Dados}
A assinatura de tipo de uma função em Haskell que recebe múltiplos argumentos é construída usando múltiplas setas (\haskellinline{->}). Considere a assinatura genérica \haskellinline{f :: Int -> Bool -> String}. 

Devido à regra de associatividade da seta de tipo em Haskell, qual é a representação matematicamente equivalente e explícita com parênteses dessa assinatura?

\begin{enumerate}
  \item \haskellinline{f :: Int -> Bool -> (String)}

  \item \haskellinline{f :: Int -> (Bool -> String)}

  \item A seta não possui associatividade; os parênteses não podem ser adicionados sem alterar fundamentalmente a aridade da função.

  \item \haskellinline{f :: (Int, Bool) -> String}

  \item \haskellinline{f :: (Int -> Bool) -> String}


\end{enumerate}

\questionTitle{11}{5: Expressão Condicional}
Considere a seguinte definição de função em Haskell, que propositalmente omite a guarda \haskellinline{otherwise}:

\begin{minted}{haskell}
classifica x
  | x > 0  = "Positivo"
  | x < 0  = "Negativo"
\end{minted}

O que acontecerá estritamente se o programa avaliar a expressão \haskellinline{classifica 0}?

\begin{enumerate}
  \item O compilador emitirá um erro de sintaxe, pois a linguagem Haskell obriga explicitamente o uso da palavra-chave \haskellinline{otherwise} no final de qualquer bloco de guardas.

  \item A função retornará silenciosamente um valor nulo ou vazio (como a string vazia \haskellinline{""}) por padrão, para evitar a interrupção do programa.

  \item Ocorrerá um erro de tipo em tempo de compilação, porque a função não consegue provar que todas as saídas pertencem ao tipo \haskellinline{String}.

  \item A avaliação entrará em um laço infinito (\emph{loop} infinito) procurando por uma condição que satisfaça a entrada.

  \item Ocorrerá um erro em tempo de execução (\emph{Non-exhaustive patterns in function}), pois nenhuma das guardas resultou em \haskellinline{True} e não há um caso de fallback definido.


\end{enumerate}

\questionTitle{12}{5: Expressão Condicional}
Analise a função a seguir. Ela faz uso de guardas e de uma cláusula \haskellinline{where}.

\begin{minted}{haskell}
analisa x
  | t > 10    = t - 10
  | otherwise = t
  where t = x * x
\end{minted}

Qual será o resultado da avaliação de \haskellinline{analisa 3} e \haskellinline{analisa 4}, respectivamente?

\begin{enumerate}
  \item \haskellinline{9} e \haskellinline{16}

  \item \haskellinline{9} e \haskellinline{6}

  \item Ocorre um erro de escopo, pois a variável \haskellinline{t} não pode ser usada na condição booleana das guardas, apenas no resultado.

  \item \haskellinline{-1} e \haskellinline{6}

  \item \haskellinline{3} e \haskellinline{4}


\end{enumerate}

\questionTitle{13}{7: Polimorfismo}
A função \haskellinline{read} converte uma \haskellinline{String} em um valor de qualquer tipo que pertença à classe \haskellinline{Read}. Sua assinatura é \haskellinline{read :: Read a => String -> a}. 
Se você digitar apenas \haskellinline{read "42"} no \emph{prompt} do GHCi e pressionar Enter, ocorrerá um erro de \textbf{ambiguidade}. Por que isso acontece?

\begin{enumerate}
  \item Porque a função \haskellinline{read} só opera sobre dados definidos pelo usuário; para tipos primitivos, deve-se usar funções específicas como \haskellinline{parseInt} ou \haskellinline{parseDouble}.

  \item Porque valores de retorno polimórficos violam as regras do paradigma funcional puro, já que o compilador não consegue garantir o tempo de execução da conversão textual.

  \item Porque a variável de tipo \haskellinline{a} está apenas no tipo de retorno, e não há contexto na expressão (como uma operação matemática ou anotação de tipo explícita) que permita ao GHCi inferir qual tipo de número instanciar.

  \item Porque o texto \haskellinline{"42"} não possui um ponto decimal, forçando o GHCi a rejeitar o polimorfismo que exigiria o formato \haskellinline{"42.0"}.

  \item Porque o GHCi não possui a classe \haskellinline{Read} importada por padrão. O aluno deve explicitamente usar \haskellinline{import Data.Read} antes de chamar a função.


\end{enumerate}

\questionTitle{14}{8: Programas Interativos}
Muitos alunos com experiência em linguagens imperativas assumem que a palavra-chave \haskellinline{return} possui o mesmo comportamento em Haskell. Analise o seguinte código interativo:

\begin{minted}{haskell}
interacao :: IO ()
interacao = do
  putStrLn "Início"
  return "Um valor qualquer"
  putStrLn "Fim"
\end{minted}

Qual afirmação descreve \textbf{corretamente} o comportamento da função \haskellinline{return} neste bloco \haskellinline{do}?

\begin{enumerate}
  \item A palavra-chave \haskellinline{return} extrai o valor puro de dentro de uma ação \haskellinline{IO}, sendo o exato oposto do operador \haskellinline{<-}.

  \item A função \haskellinline{return} encerra a execução do bloco \haskellinline{do} imediatamente, fazendo com que a mensagem \haskellinline{"Fim"} nunca seja exibida na tela.

  \item Ocorre um erro de tipo, pois \haskellinline{return} tenta alterar a assinatura da função de \haskellinline{IO ()} para \haskellinline{IO String}, o que é proibido no meio de um bloco.

  \item A função \haskellinline{return} atua como um comando de saída padrão, imprimindo \haskellinline{"Um valor qualquer"} no terminal do usuário.

  \item A função \haskellinline{return} apenas empacota a string em uma ação do tipo \haskellinline{IO String}, mas não interrompe o fluxo de execução; a mensagem \haskellinline{"Fim"} será impressa normalmente na tela.


\end{enumerate}

\questionTitle{15}{3: Expressões e Definições}
Uma das regras mais importantes de precedência em Haskell diz respeito à aplicação de função. Dado que \haskellinline{length} retorna um inteiro (\haskellinline{Int}) indicando o tamanho de uma lista, o que acontecerá ao tentar avaliar a expressão abaixo?
\begin{minted}{haskell}
length [1, 2, 3] ++ [4, 5]
\end{minted}

\begin{enumerate}
  \item O resultado será \haskellinline{8}. A função \haskellinline{length} é distribuída automaticamente para todos os operandos do operador \haskellinline{++}, somando os tamanhos de ambas as listas.

  \item Ocorrerá um erro de tipo. A aplicação de função prefixa tem a maior precedência de todas, avaliando \haskellinline{length [1, 2, 3]} primeiro, e em seguida tentando aplicar o operador de concatenação (\haskellinline{++}) a um número inteiro e uma lista.

  \item Ocorrerá um erro de sintaxe. O uso de operadores infixos requer que todos os argumentos ao seu redor estejam explicitamente agrupados com parênteses, como em \haskellinline{(length [1, 2, 3]) ++ [4, 5]}.

  \item O resultado será \haskellinline{5}. O operador infixo \haskellinline{++} tem maior precedência e concatenará as listas primeiro, resultando em \haskellinline{length [1, 2, 3, 4, 5]}.

  \item O resultado será \haskellinline{[3, 4, 5]}. O Haskell avaliará \haskellinline{length [1, 2, 3]} como \haskellinline{3}, e converterá esse valor em uma lista implicitamente para a concatenação.


\end{enumerate}

\questionTitle{16}{10: Ações de E/S Recursivas}
Sabemos que a avaliação preguiçosa (\emph{lazy evaluation}) permite processar listas puras infinitas (ex: \haskellinline{take 3 (geraSequencia 0)}). 
Se você criar uma função \haskellinline{leInfinito :: IO [String]} que usa recursão sem caso base para ler repetidamente com \haskellinline{getLine}, será possível extrair os 3 primeiros elementos digitados com segurança (ex: passando o resultado para uma função de limite)?

\begin{enumerate}
  \item Sim. Graças ao motor preguiçoso do GHC, a função de E/S será executada apenas três vezes, solicitando exatamente três entradas do usuário antes de encerrar o programa.

  \item Não, porque \haskellinline{leInfinito} retornará uma lista do tipo \haskellinline{[IO String]}, que é incompatível com a função padrão \haskellinline{take}, que exige uma lista pura \haskellinline{[a]}.

  \item Não. As ações na mônada \haskellinline{IO} impõem um sequenciamento estrito. O programa ficará travado em um laço infinito solicitando entrada no terminal, impedindo a geração final da lista para avaliação preguiçosa.

  \item Não, pois funções recursivas sem caso base explícito são estaticamente rejeitadas e barradas pelo compilador GHC quando contêm comandos de E/S.

  \item Sim, mas o programa exibirá um \emph{warning} no terminal alertando que o fluxo de entrada (\emph{stdin}) não foi fechado corretamente, embora os 3 itens sejam avaliados.


\end{enumerate}

\questionTitle{17}{1: Paradigmas}
Em Haskell, a execução de um programa ocorre pela avaliação de expressões através da \textbf{redução} (substituição de parâmetros). Considere a expressão \haskellinline{dobro (dobro 2)}, onde \haskellinline{dobro x = x + x}. 

Sobre as estratégias de avaliação \emph{de dentro para fora} e \emph{de fora para dentro}, qual afirmação é estritamente correta segundo os princípios do paradigma puramente funcional?

\begin{enumerate}
  \item Ambas as estratégias alteram a transparência referencial da expressão durante os passos intermediários, mas o resultado final converge.

  \item A ordem em que as reduções são realizadas não afeta o resultado final da expressão, uma propriedade garantida pela ausência de efeitos colaterais.

  \item O compilador Haskell proíbe a avaliação \emph{de fora para dentro} porque ela duplica a expressão interna, causando vazamento de memória.

  \item A avaliação \emph{de fora para dentro} é a única permitida pela linguagem, pois garante que o estado global não seja alterado prematuramente.

  \item A avaliação \emph{de dentro para fora} pode resultar em um valor diferente se a função \haskellinline{dobro} for aplicada a números negativos.


\end{enumerate}

\questionTitle{18}{1: Paradigmas}
Um compilador Haskell pode realizar uma otimização agressiva chamada \emph{memoização} (cache), substituindo todas as chamadas subsequentes de \haskellinline{calculaRisco 5000} pelo resultado pré-calculado da primeira chamada. 

Qual propriedade fundamental da linguagem garante que essa substituição é \textbf{matematicamente segura} e não introduzirá falhas no sistema?

\begin{enumerate}
  \item O sistema de inferência de tipos Hindley-Milner, que prova estaticamente que a função não pode gerar exceções de divisão por zero.

  \item A transparência referencial, que assegura que uma função pura sempre retornará o mesmo resultado para a mesma entrada, sem depender de contexto ou estado externo.

  \item A tipagem estática forte, que impede que variáveis globais sejam alteradas por \textit{threads} concorrentes durante o cálculo.

  \item A avaliação preguiçosa (\emph{lazy evaluation}), que garante que \haskellinline{calculaRisco} só será efetivamente executada quando o usuário imprimir o resultado na tela.

  \item O polimorfismo paramétrico, que permite que o resultado seja armazenado em cache independentemente de o argumento ser \haskellinline{Int} ou \haskellinline{Double}.


\end{enumerate}

\questionTitle{19}{4: Tipos de Dados}
Haskell oferece dois tipos inteiros básicos muito utilizados: \haskellinline{Int} e \haskellinline{Integer}. Se você precisar calcular e armazenar o fatorial de 1000 ($1000!$, um número com mais de 2500 dígitos) em um programa Haskell, qual o tipo mais apropriado a se utilizar na assinatura da função e por quê?

\begin{enumerate}
  \item Ambos falharão. Haskell exige o uso de bibliotecas de terceiros como \haskellinline{BigMath} para cálculos que superam os limites do hardware de 64 bits.

  \item \haskellinline{Integer}, pois é um tipo de precisão arbitrária que pode representar números inteiros de qualquer tamanho, limitado apenas pela memória RAM disponível.

  \item \haskellinline{Double}, pois a notação científica de ponto flutuante é a única forma suportada pela linguagem para processar e exibir números com mais de 2500 dígitos.

  \item \haskellinline{Int}, pois, diferentemente de linguagens antigas como C, o \haskellinline{Int} em Haskell escala seu tamanho automaticamente na memória para evitar \emph{overflow}.

  \item \haskellinline{Integer}, pois é o único tipo em Haskell que aceita operações recursivas profundas; o \haskellinline{Int} é estritamente usado para contadores de repetição simples.


\end{enumerate}

\questionTitle{20}{7: Polimorfismo}
Qual é a diferença fundamental e arquitetural entre o comportamento de uma função com \textbf{polimorfismo paramétrico} (ex: \haskellinline{length :: [a] -> Int}) e uma com \textbf{polimorfismo ad hoc} (ex: \haskellinline{sum :: Num a => [a] -> a})?

\begin{enumerate}
  \item A função paramétrica opera exatamente com a mesma lógica e código subjacente independentemente do tipo instanciado, enquanto a função \emph{ad hoc} pode ter implementações completamente diferentes dependendo do tipo específico usado.

  \item Não há diferença no comportamento em tempo de execução; a distinção é puramente sintática para organizar o código-fonte em módulos diferentes.

  \item Funções paramétricas violam a transparência referencial porque o tipo de retorno pode variar, enquanto funções \emph{ad hoc} são sempre estritamente puras.

  \item A função paramétrica só pode receber tipos primitivos (como \haskellinline{Int} ou \haskellinline{Char}), enquanto o polimorfismo \emph{ad hoc} permite o uso de estruturas de dados complexas criadas pelo usuário.

  \item O polimorfismo paramétrico exige que o programador declare o tipo explicitamente na chamada da função, enquanto o \emph{ad hoc} usa a inferência de tipos do compilador GHC para descobrir o tipo em tempo de execução.


\end{enumerate}

\questionTitle{21}{1: Paradigmas}
A \emph{crise do software} é caracterizada por desafios crescentes de custo, complexidade e, criticamente, falta de garantia de correção em sistemas vitais. De acordo com o texto, como a adoção do paradigma funcional ataca diretamente o problema da \textbf{confiança e correção} do software?

\begin{enumerate}
  \item Por obrigar o uso de estruturas de dados orientadas a objetos, o paradigma garante que os estados inválidos sejam rejeitados logo na instanciação das classes.

  \item Por compilar o código diretamente para linguagem de máquina otimizada, eliminando a camada de máquina virtual que frequentemente introduz vulnerabilidades de segurança.

  \item Por utilizar a manipulação explícita de ponteiros em ambientes isolados (\emph{sandboxes}), evitando que o estado global seja corrompido.

  \item Por tratar o programa como expressões matemáticas sem efeitos colaterais, o paradigma incentiva e facilita a prova formal de que o código se comporta conforme o especificado.

  \item Por substituir a necessidade de testes de unidade automatizados pela inferência de tipos, que detecta 100\% dos erros lógicos em tempo de compilação.


\end{enumerate}

\questionTitle{22}{5: Expressão Condicional}
Como \haskellinline{if-then-else} em Haskell é uma expressão (e não uma declaração de controle de fluxo), ele pode resultar em qualquer tipo de valor, inclusive em outras funções. 

Qual será o resultado da avaliação da seguinte expressão?
\begin{minted}{haskell}
(if 5 > 3 then (*) else (+)) 4 2
\end{minted}

\begin{enumerate}
  \item \haskellinline{14}

  \item Ocorre um erro de tipo, pois os ramos \haskellinline{then} e \haskellinline{else} de um \haskellinline{if} não podem conter operadores puros ou funções.

  \item \haskellinline{6}

  \item Ocorre um erro de sintaxe, pois não é permitido colocar parênteses ao redor de um \haskellinline{if-then-else}.

  \item \haskellinline{8}


\end{enumerate}

\questionTitle{23}{8: Programas Interativos}
Considere a seguinte ação interativa, onde um aluno usou \haskellinline{let} em vez do operador \haskellinline{<-} por engano:

\begin{minted}{haskell}
saudacao :: IO ()
saudacao = do
  putStrLn "Qual o seu nome?"
  let nome = getLine
  putStrLn "Olá!"
\end{minted}

O que acontecerá quando a ação \haskellinline{saudacao} for avaliada e executada pelo sistema de \emph{runtime}?

\begin{enumerate}
  \item O programa entrará em um laço infinito aguardando a entrada do usuário, pois o \haskellinline{let} desabilita o \emph{buffering} padrão da entrada do sistema.

  \item A função \haskellinline{getLine} lançará uma exceção em tempo de execução porque o compilador tentará converter um tipo de E/S em um tipo puro à força.

  \item O programa não aguardará a digitação do usuário. A palavra \haskellinline{let} apenas associa a variável \haskellinline{nome} à ação (uma \emph{receita} do tipo \haskellinline{IO String}), mas não a executa.

  \item Ocorre um erro de sintaxe imediato, pois a palavra-chave \haskellinline{let} é estritamente proibida dentro de blocos \haskellinline{do}.

  \item O programa aguardará a entrada do usuário normalmente e salvará a string lida na variável \haskellinline{nome}, mas ocorrerá um aviso (\emph{warning}) de estilo.


\end{enumerate}

\questionTitle{24}{8: Programas Interativos}
No trecho de código a seguir, a ação \haskellinline{getLine} foi colocada isoladamente em uma linha do bloco \haskellinline{do}, sem o uso do operador de vinculação (\haskellinline{<-}):

\begin{minted}{haskell}
pausa :: IO ()
pausa = do
  putStrLn "Pressione Enter para continuar..."
  getLine
  putStrLn "Continuando."
\end{minted}

Qual será o resultado pragmático dessa implementação durante a execução do programa?

\begin{enumerate}
  \item A string capturada pelo \haskellinline{getLine} será implicitamente passada como argumento para a função \haskellinline{putStrLn} da linha seguinte, imprimindo o que o usuário digitou.

  \item Ocorrerá um erro de compilação (\emph{Parse error in pattern}), visto que o Haskell obriga o armazenamento de qualquer valor gerado em um bloco \haskellinline{do}.

  \item O compilador ignorará completamente a linha do \haskellinline{getLine} como código morto (\emph{dead code}), pois seu resultado não foi atribuído a nenhuma variável.

  \item A ação \haskellinline{getLine} será efetivamente executada, pausando o programa até que o usuário pressione Enter, mas a string digitada será descartada e não poderá ser usada posteriormente.

  \item O programa saltará o \haskellinline{getLine} se não houver um teclado físico disponível, retornando imediatamente \haskellinline{Nothing}.


\end{enumerate}

\questionTitle{25}{10: Ações de E/S Recursivas}
Considere o fragmento de código abaixo, que lê uma lista de strings do teclado de forma recursiva até encontrar uma linha vazia:

\begin{minted}{haskell}
leTextos :: IO [String]
leTextos = do
  t <- getLine
  if t == "" 
    then return []
    else do
      ts <- leTextos
      return (t:ts)
\end{minted}

No ramo \haskellinline{else}, qual é o papel estrutural exato e inegociável da instrução \haskellinline{return (t:ts)}?

\begin{enumerate}
  \item A instrução \haskellinline{return} serve apenas para indicar ao compilador que a recursão terminou; a construção real da lista na memória já ocorreu implicitamente na linha \haskellinline{ts <- leTextos}.

  \item Ela extrai a string armazenada na variável \haskellinline{t} de dentro do contexto \haskellinline{IO}, permitindo que seja anexada à lista pura \haskellinline{ts}.

  \item Ela recebe a lista pura construída pelo operador \haskellinline{(:)} e a empacota de volta em um contexto \haskellinline{IO [String]}, satisfazendo a assinatura da ação como o último passo do bloco \haskellinline{do}.

  \item Ela instrui o ambiente de execução a retornar o controle para o sistema operacional, finalizando as chamadas recursivas empilhadas e liberando a memória do terminal.

  \item Trata-se de uma redundância permitida sintaticamente, uma vez que omitir o \haskellinline{return} e escrever apenas \haskellinline{t:ts} também resultaria na mesma compilação válida.


\end{enumerate}

\questionTitle{26}{10: Ações de E/S Recursivas}
Um aluno tentou implementar uma contagem regressiva interativa, mas cometeu um erro sintático induzido por sua vivência em linguagem C:

\begin{minted}{haskell}
contagem :: Int -> IO ()
contagem 0 = putStrLn "Fogo!"
contagem n = do
  putStrLn (show n)
  return (contagem (n - 1))
\end{minted}

O compilador emitirá um erro de incompatibilidade de tipos (\emph{Type mismatch}). Qual é a justificativa rigorosa para isso e como consertar?

\begin{enumerate}
  \item A operação matemática \haskellinline{n - 1} quebra a inferência de tipo dentro da diretiva \haskellinline{return}. A solução é realizar o cálculo fora: \haskellinline{let prox = n - 1 in return (contagem prox)}.
  \item A função \haskellinline{return} exige como parâmetro um valor primitivo (\haskellinline{Int}, \haskellinline{Bool}), não podendo receber a aplicação de uma função recursiva. A correção é usar um \haskellinline{let x = contagem (n-1)}.

  \item Em blocos \haskellinline{do}, a palavra \haskellinline{return} só pode aparecer como o primeiro comando, servindo de ponto de inicialização. O compilador acusa erro por estar na última linha.

  \item O erro está no uso de parênteses dentro do \haskellinline{return}, forçando a criação de uma tupla monádica. A solução é remover os parênteses externos: \haskellinline{return contagem (n - 1)}.

  \item A função \haskellinline{contagem} já resulta em \haskellinline{IO ()}. Ao usar \haskellinline{return}, o aluno cria o tipo \haskellinline{IO (IO ())}. A correção é simplesmente chamar a recursão: escrever \haskellinline{contagem (n - 1)} como a última linha do bloco.


\end{enumerate}

\questionTitle{27}{9: Recursividade}
Um aluno implementou a seguinte função para contar quantos números pares existem em uma lista:

\begin{minted}{haskell}
contaPares :: [Int] -> Int
contaPares (x:xs)
  | even x    = 1 + contaPares xs
  | otherwise = contaPares xs
\end{minted}

O que acontecerá estritamente se o aluno avaliar a expressão \haskellinline{contaPares [2, 4, 6]} no GHCi?

\begin{enumerate}
  \item A função processará todos os elementos, mas ao tentar avaliar o resto vazio no final da lista, lançará uma exceção em tempo de execução (\emph{Non-exhaustive patterns}) porque o caso base \haskellinline{contaPares []} foi omitido.

  \item A função retornará \haskellinline{1}, pois o operador de adição não suporta acumulação recursiva; ele avaliará apenas o primeiro elemento par encontrado.

  \item A expressão será avaliada perfeitamente e retornará \haskellinline{3}, pois o compilador GHC infere automaticamente que listas vazias devem retornar \haskellinline{0} em contextos numéricos.

  \item Ocorrerá um erro em tempo de compilação (\emph{Parse error}), pois a linguagem proíbe o uso de guardas dentro de definições de funções recursivas.

  \item O programa entrará em um laço infinito, pois não há uma instrução de parada explícita como um \haskellinline{return} ou um \haskellinline{break}.


\end{enumerate}

\questionTitle{28}{9: Recursividade}
Considere a seguinte função recursiva misteriosa implementada em Haskell:

\begin{minted}{haskell}
misterio :: [a] -> [a]
misterio [] = []
misterio (x:xs) = misterio xs ++ [x]
\end{minted}

Avaliando mentalmente os passos de redução (empilhamento de chamadas), qual será o resultado exato da expressão \haskellinline{misterio [1, 2, 3]}?

\begin{enumerate}
  \item \haskellinline{[1, 2, 3]}. A função é equivalente à função identidade, recriando a lista original na mesma ordem.

  \item Haverá um aviso de compilação indicando que a variável de tipo \haskellinline{a} não suporta a operação de soma implícita na construção.

  \item Ocorrerá um erro de tipo, pois não é possível usar o operador de concatenação (\haskellinline{++}) para juntar o resultado de uma função com um literal de lista \haskellinline{[x]}.

  \item \haskellinline{[3, 2, 1]}. A função inverte a lista ao colocar o primeiro elemento no final da lista resultante das chamadas recursivas sucessivas.

  \item \haskellinline{[[1], [2], [3]]}. O operador de construção transforma cada elemento individual em uma sublista antes de concatenar.


\end{enumerate}

\questionTitle{29}{2: Ambiente}
Durante a exploração de bibliotecas no ambiente interativo (GHCi), você deseja alcançar dois objetivos em sequência:
Primeiro, listar \textbf{todas} as funções disponíveis dentro do módulo \haskellinline{Data.List}.
Segundo, descobrir a \textbf{assinatura de tipo} da função \haskellinline{delete}, sem executá-la.

Quais comandos internos do GHCi você deve utilizar para realizar essas duas tarefas, respectivamente?

\begin{enumerate}
  \item \texttt{:browse Data.List} e, em seguida, \texttt{:type delete} (ou \texttt{:t delete}).

  \item \texttt{:show Data.List} e, em seguida, \texttt{:eval delete}.

  \item \texttt{:load Data.List} e, em seguida, \texttt{:help delete}.

  \item \texttt{:type Data.List} e, em seguida, \texttt{:browse delete}.

  \item \texttt{:info Data.List} e, em seguida, \texttt{:browse delete}.


\end{enumerate}

\questionTitle{30}{7: Polimorfismo}
Considere a seguinte definição de função, que compara dois valores polimórficos:
\begin{minted}{haskell}
analisa x y = (x == y, x < y)
\end{minted}
O operador \haskellinline{==} exige a restrição \haskellinline{Eq a =>}, enquanto o operador \haskellinline{<} exige a restrição \haskellinline{Ord a =>}. Qual será a assinatura de tipo inferida pelo compilador GHC para a função \haskellinline{analisa}?

\begin{enumerate}
  \item \haskellinline{analisa :: a -> a -> (Bool, Bool)}. Por serem operadores primitivos da linguagem, o compilador suspende a checagem de classes de tipo e os assume como totalmente polimórficos.

  \item \haskellinline{analisa :: (Eq a, Ord a) => a -> a -> (Bool, Bool)}. O compilador sempre lista detalhadamente todas as restrições de operadores utilizados na função, sem aplicar hierarquias lógicas.

  \item Ocorrerá um erro de inferência de tipo, pois uma função não pode retornar uma tupla gerada por operadores que pertencem a classes de tipo diferentes sem uma conversão explícita.

  \item \haskellinline{analisa :: (Ord a, Ord b) => a -> b -> (Bool, Bool)}. Como os valores podem ser diferentes (ex: \haskellinline{5} e \haskellinline{5.0}), o compilador infere variáveis de tipo independentes para cada argumento.

  \item \haskellinline{analisa :: Ord a => a -> a -> (Bool, Bool)}. A restrição \haskellinline{Eq a} é omitida porque, na hierarquia de classes, ser uma instância de \haskellinline{Ord} obrigatoriamente implica ser uma instância de \haskellinline{Eq}.


\end{enumerate}

\questionTitle{31}{6: Estruturas de dados básicas}
As funções \haskellinline{fst} e \haskellinline{snd} são fornecidas pelo Prelude para acessar elementos de uma tupla. O que acontecerá se tentarmos avaliar a expressão \haskellinline{fst (10, 20, 30)} no GHCi?

\begin{enumerate}
  \item Ocorrerá um erro de execução, pois o GHCi tentará extrair o primeiro elemento, mas se perderá ao encontrar o terceiro valor na memória.

  \item O resultado será \haskellinline{(10, 20)}, pois em tuplas maiores que 2 elementos, \haskellinline{fst} retorna os dois primeiros valores por padrão.

  \item O resultado será \haskellinline{10}, pois \haskellinline{fst} sempre extrai o primeiro elemento de qualquer tupla, independentemente de sua aridade.

  \item A expressão será avaliada corretamente se importarmos o módulo \haskellinline{Data.Tuple} antes da execução.

  \item Ocorrerá um erro de tipo em tempo de compilação. As funções \haskellinline{fst} e \haskellinline{snd} são definidas estritamente para pares (tuplas de dois elementos), e não funcionam em triplas.


\end{enumerate}

\questionTitle{32}{4: Tipos de Dados}
Linguagens com tipagem dinâmica (como Python ou JavaScript) permitem criar listas mistas, como \texttt{[1, True, "texto"]}. Em Haskell, se um programador tentar avaliar a expressão \haskellinline{[1, True, "texto"]}, o que acontecerá e por quê?

\begin{enumerate}
  \item Ocorrerá um erro apenas em tempo de execução, caso o programa tente realizar uma operação matemática (como soma) sobre o elemento booleano ou textual da lista.

  \item O compilador emitirá um aviso (\emph{warning}), mas criará uma tupla automaticamente para acomodar os tipos diferentes.

  \item A expressão será aceita, mas o tipo inferido será o tipo universal \haskellinline{Any}, desativando as verificações de tipo estáticas para essa lista.

  \item Ocorrerá um erro em tempo de compilação. Devido à tipagem estática e forte de Haskell, os elementos de uma lista devem pertencer estritamente ao mesmo tipo (homogeneidade).

  \item A expressão será convertida implicitamente, transformando todos os valores em \haskellinline{String} (resultando em \haskellinline{["1", "True", "texto"]}) para manter a compatibilidade.


\end{enumerate}

\questionTitle{33}{8: Programas Interativos}
Um programador deseja criar uma função puramente matemática que calcule o quadrado de um número, mas que também imprima o resultado na tela como um efeito colateral de \emph{log} (registro). Ele escreve a seguinte tentativa:

\begin{minted}{haskell}
quadrado :: Int -> Int
quadrado x = 
  let res = x * x
  in putStrLn ("Log: " ++ show res)
\end{minted}

Por que o compilador GHC rejeitará este código terminantemente?

\begin{enumerate}
  \item Porque a sintaxe do \haskellinline{let...in} não suporta a chamada de funções que recebam \haskellinline{String} como parâmetro.

  \item Porque o sistema de tipos garante a pureza: a função \haskellinline{putStrLn} retorna uma ação \haskellinline{IO ()}, que não pode ser devolvida como resultado de uma função que promete retornar um valor puro \haskellinline{Int}.

  \item O compilador não rejeitará o código; ele irá compilá-lo, mas o \emph{log} será suprimido na tela para garantir a pureza da função em tempo de execução.

  \item Porque a função \haskellinline{show} gera um conflito de precedência com a concatenação (\haskellinline{++}), exigindo parênteses adicionais.

  \item Porque todas as funções em Haskell são obrigadas a usar a anotação do bloco \haskellinline{do}, independentemente de possuírem efeitos colaterais ou não.


\end{enumerate}

\questionTitle{34}{4: Tipos de Dados}
Considere as seguintes declarações de tipos sinônimos e de uma função em um módulo Haskell:

\begin{minted}{haskell}
type Nome = String
type Endereco = String

imprimeEtiqueta :: Nome -> Endereco -> IO ()
imprimeEtiqueta n e =
  putStrLn ("Para: " ++ n ++ " - Env: " ++ e)
\end{minted}

Se um programador acidentalmente inverter a ordem dos argumentos e chamar \haskellinline{imprimeEtiqueta "Rua X" "João"}, o que o compilador GHC fará?

\begin{enumerate}
  \item O compilador rejeitará o código com um erro de tipo, protegendo o sistema contra a inversão de parâmetros de domínios diferentes. Esta é a principal vantagem dos tipos sinônimos.

  \item O código compilará, mas lançará uma exceção em tempo de execução porque o sistema de \emph{runtime} não conseguirá realizar o \emph{cast} de \haskellinline{Endereco} para \haskellinline{Nome}.

  \item O compilador emitirá um aviso (\emph{warning}) sobre possível mistura de tipos sinônimos, mas gerará o executável normalmente.

  \item Ocorrerá um erro sintático, pois a declaração de tipos sinônimos obriga o programador a instanciar os valores usando construtores (ex: \haskellinline{Nome "João"}).

  \item O código compilará perfeitamente sem nenhum erro ou aviso. Tipos sinônimos criados com a palavra-chave \haskellinline{type} são apenas apelidos (\emph{aliases}); para o compilador, \haskellinline{Nome}, \haskellinline{Endereco} e \haskellinline{String} são o mesmo tipo.


\end{enumerate}

\questionTitle{35}{3: Expressões e Definições}
As expressões \haskellinline{let...in} podem ser aninhadas livremente em Haskell, criando escopos locais. Qual será o resultado final da avaliação da expressão a seguir?
\begin{minted}{haskell}
let x = 10 
in (let x = 2; y = x * 3 in x + y) * x
\end{minted}

\begin{enumerate}
  \item Ocorre um erro de escopo porque a variável \haskellinline{x} é declarada duas vezes na mesma expressão, violando a regra de unicidade de variáveis.

  \item \haskellinline{8}

  \item \haskellinline{16}

  \item \haskellinline{320}

  \item \haskellinline{80}


\end{enumerate}

\questionTitle{36}{2: Ambiente}
O ecossistema Haskell divide as bibliotecas em três níveis práticos de disponibilidade. Suponha que um desenvolvedor queira utilizar:
1. A função \haskellinline{sqrt} (raiz quadrada).
2. A função \haskellinline{delete} do módulo \haskellinline{Data.List}.
3. Uma função de uma biblioteca de rede chamada \haskellinline{network}, desenvolvida por terceiros.

Como o desenvolvedor deve acessar essas três funcionalidades, respectivamente?

\begin{enumerate}
  \item 1. Usar diretamente; 2. Instalar o pacote \texttt{base} via Cabal e importá-lo; 3. Solicitar que o GHCup compile a ferramenta de rede em tempo de execução.

  \item 1. Importar o módulo \texttt{Prelude}; 2. Usar diretamente (é carregada automaticamente); 3. Baixar o arquivo \texttt{.hs} do Hackage manualmente e colocá-lo na pasta do projeto.

  \item Todas as três requerem a instalação prévia via Cabal, pois o GHC, por padrão, fornece apenas o compilador sem nenhuma biblioteca embutida.

  \item 1. Usar diretamente (pertence ao \texttt{Prelude}); 2. Importar o módulo explicitamente (vem com a Biblioteca Padrão do GHC); 3. Instalar o pacote via Cabal a partir do Hackage e depois importá-lo.

  \item 1. Usar diretamente; 2. Importar o módulo explicitamente; 3. Usar diretamente também, pois o GHC moderno carrega o Hackage de forma transparente (\emph{lazy loading}).


\end{enumerate}

\questionTitle{37}{6: Estruturas de dados básicas}
Considere uma função hipotética \haskellinline{buscaIdade :: String -> Maybe Int}, que busca a idade de um usuário pelo nome e retorna \haskellinline{Nothing} se não encontrar. 

Um aluno tentou escrever a seguinte expressão para calcular o ano de nascimento aproximado (considerando o ano atual como 2026):
\begin{minted}{haskell}
2026 - buscaIdade "Alice"
\end{minted}

Qual o problema com esta expressão sob a ótica do sistema de tipos de Haskell?

\begin{enumerate}
  \item Nenhum problema. Se \haskellinline{buscaIdade} retornar \haskellinline{Nothing}, a expressão inteira resultará em \haskellinline{0}, avaliando para \haskellinline{2026 - 0 = 2026}.

  \item O erro ocorre porque a função \haskellinline{buscaIdade} deveria retornar uma string vazia (\haskellinline{""}) em vez de \haskellinline{Nothing} para ser compatível com as regras de E/S.

  \item Haverá um laço infinito, pois o GHCi tentará extrair o valor de \haskellinline{Nothing} repetidamente caso a usuária não seja encontrada.

  \item A subtração deveria ser invertida (\haskellinline{buscaIdade "Alice" - 2026}), pois valores envelopados em construtores como \haskellinline{Maybe} só podem aparecer no lado esquerdo de operadores infixos.

  \item O operador de subtração (\haskellinline{-}) exige tipos numéricos, mas \haskellinline{buscaIdade} retorna um \haskellinline{Maybe Int}. Não é possível realizar operações aritméticas diretamente sobre o construtor \haskellinline{Just} ou \haskellinline{Nothing} sem antes desempacotar o valor.


\end{enumerate}

\questionTitle{38}{6: Estruturas de dados básicas}
Sabendo que em Haskell o tipo \haskellinline{String} é apenas um sinônimo para \haskellinline{[Char]} (uma lista de caracteres), qual será o resultado da avaliação da seguinte expressão?
\begin{minted}{haskell}
['H', 'a'] ++ "skell"
\end{minted}

\begin{enumerate}
  \item Ocorrerá um erro de tipo, pois o operador \haskellinline{++} exige que ambos os lados tenham exatamente a mesma sintaxe literal (ou ambos com aspas duplas, ou ambos com aspas simples).

  \item O resultado será uma tupla mista contendo a lista e a string: \haskellinline{(['H', 'a'], "skell")}.

  \item O resultado será uma lista de strings: \haskellinline{["Ha", "skell"]}.

  \item Ocorrerá um erro de sintaxe, pois listas de caracteres definidas com aspas simples não podem ser concatenadas diretamente sem o uso da função \haskellinline{show}.

  \item \haskellinline{"Haskell"} (ou \haskellinline{['H', 'a', 's', 'k', 'e', 'l', 'l']}). Como ambos os operandos são listas de caracteres, o operador de concatenação (\haskellinline{++}) junta as duas listas perfeitamente.


\end{enumerate}

\questionTitle{39}{9: Recursividade}
A recursão de cauda (\emph{tail recursion}) é uma técnica vital em programação funcional onde a chamada recursiva é a \textbf{última} operação a ser avaliada, permitindo que o compilador otimize a execução e evite o estouro da pilha de chamadas. 

Qual das alternativas abaixo apresenta uma implementação com recursão de cauda correta?

\begin{enumerate}
  \item \begin{minted}{haskell}
somaNormal [] = 0
somaNormal (x:xs) = x + somaNormal xs
\end{minted}

  \item \begin{minted}{haskell}
somaCauda acc [] = acc
somaCauda acc (x:xs) = somaCauda (acc + x) xs
\end{minted}

  \item \begin{minted}{haskell}
tamanho [] = 0
tamanho (_:xs) = 1 + tamanho xs
\end{minted}

  \item \begin{minted}{haskell}
somaErro (x:xs) = 
  let resultado = somaErro xs 
  in x + resultado
\end{minted}

  \item \begin{minted}{haskell}
dobra [] = []
dobra (x:xs) = (x * 2) : dobra xs
\end{minted}


\end{enumerate}

\questionTitle{40}{5: Expressão Condicional}
Embora o uso de guardas seja recomendado para melhor legibilidade, é perfeitamente legal aninhar expressões \haskellinline{if-then-else} em Haskell. 

Avalie a função abaixo com a entrada \haskellinline{f 5}. Qual será o resultado gerado?
\begin{minted}{haskell}
f x = if x > 0
      then if x > 10
           then 1
           else 2
      else 3
\end{minted}

\begin{enumerate}
  \item Ocorre um erro de sintaxe, pois faltam parênteses obrigatórios no aninhamento.

  \item \haskellinline{1}

  \item \haskellinline{3}

  \item \haskellinline{2}

  \item Ocorre um erro de ambiguidade; o compilador não sabe a qual \haskellinline{if} o primeiro \haskellinline{else} pertence.


\end{enumerate}


\end{multicols}
\makeanswersheet{UFOP}{BCC222: Programação Funcional}{José Romildo}{2025/2}{Exame Especial Parcial 1}{03}{40}{25}
\gabarito{03}{B,C,C,C,C,D,E,B,C,B,E,B,C,E,B,C,B,B,B,A,D,E,C,D,C,E,A,D,A,E,E,D,B,E,E,D,E,E,B,D}
\end{document}
